\scanpage{57}
\mspage{28}
would therefore group together \inserted{\unsure{classes of|\deleted{\unsure{classes} of}}} \unsure{\deleted{\unsure{in}\ \unsure{some}}|in some} packets
(by looking at the centralizers of the image subgroups
of $\Gamma^\circ$ under its sections,\deleted{\unsure{)|f|and}} and \unsure{by|a}
putting together splittings which have conjugate
$\unsure{\mathrm{Zentr}|\mathrm{Centr}}(\Gamma^\circ)$ (even if they are \unsure{not|thereof}
themselves conjugate). Thus one again recovers
a description of $X(K)$ (thus
also of the $X(K')$ for every \unsure{finite|finite} extension \inserted{$K'$} of $K$)
in terms of the extension $E_L$ of $E_K$ by
$\Pi_{L/K}$, at the $\unsure{=|\overline m}$ time as a way of
reconstructing the $U=X-\overline{\unsure{I|\underline I}}$\ldots
Thus in fact it is the structure $E_L\unsure{-|\longrightarrow}E_K$
which is the richest a priori, and far
more convenient for genus $0$ and $1$,
where the consideration of the $U$ of type $(g,\nu)$\unsure{.|,}
$(2g\unsure{+}\nu\geqslant3)$ breaks the ``continuous'' group
of automorphisms\ldots The ``model'' form
of the preceding conjecture comes down
to the ``birational'' form, when one
makes the latter precise there by saying that every
\margin{56}
